
\chapter{La tirannia delle infinite possibilità}
\label{cap:p2c21}

\epigraph{«L'uomo è condannato a essere libero.»}{Jean-Paul Sartre\index[main]{Sartre, Jean-Paul}, \textit{L'esistenzialismo è un umanismo} (1946)}

Ogni sera, milioni di persone aprono un catalogo di streaming, uno store di app o una libreria musicale con più titoli di quanti potranno mai guardare. Poi scrollano, paralizzate dall'abbondanza, e vanno a dormire senza aver visto nulla.

Eppure, se chiedo se preferirebbero un negozio sovietico degli anni '70 (un tipo di pane, uno di marmellata, uno di sapone) o un supermercato americano con 300 tipi di cereali e 40 dentifrici, scelgono il secondo. Chi vorrebbe tornare alla scarsità forzata? Ma davanti a un catalogo infinito, la libertà di scelta si trasforma nella sua nemesi.

\section{Il supermercato e il cervello che non sa fermarsi}

Sheena Iyengar e Mark Lepper condussero nel 2000 quello che sarebbe diventato uno degli esperimenti più citati della psicologia del consumo\footnote{Iyengar, S. S., \& Lepper, M. R. (2000). «When Choice is Demotivating: Can One Desire Too Much of a Good Thing?». \textit{Journal of Personality and Social Psychology}, 79(6), 995-1006.}. In un supermercato di lusso della California allestirono un banco di assaggio di marmellate: alcuni giorni presentava 24 varietà, altri solo 6. Il banco con 24 marmellate attirava più curiosi (60\% contro 40\%), ma solo il 3\% di questi acquistava; con 6 opzioni, comprava il 30\%. Più scelta, meno azione. E il dettaglio più rivelatore arrivava dopo: chi aveva comprato dal display limitato si diceva più soddisfatto settimane più tardi, mentre chi aveva scelto tra 24 opzioni continuava a chiedersi se avesse fatto la scelta giusta.

Per capire perché, dobbiamo tornare a quella stratificazione evolutiva del cervello di cui abbiamo parlato. Il nostro sistema decisionale si è formato su scelte immediate e poco numerose: combatti o fuggi, mangia o non mangiare, fidati o diffida. Anche le decisioni più complesse dei nostri antenati, quale territorio esplorare, con chi allearsi, coinvolgevano una manciata di opzioni concrete. Oggi apriamo Spotify e troviamo milioni di brani, apriamo Amazon e troviamo milioni di prodotti, e il cervello paleolitico applica a questa valanga i suoi algoritmi ancestrali. Il risultato è un cortocircuito.

Scegliere, del resto, ha un costo: dopo una serie di decisioni l'autocontrollo tende a calare\footnote{Vohs, K. D., et al. (2008). «Making choices impairs subsequent self-control: A limited-resource account of decision making, self-regulation, and active initiative». \textit{Journal of Personality and Social Psychology}, 94(5), 883-898.}. È plausibile che, dopo aver passato in rassegna molte opzioni, anche la nostra capacità decisionale si affatichi. È quello che i ricercatori chiamano «decision fatigue», affaticamento decisionale.

\section{Il paradosso digitale della scelta}

La rivoluzione digitale ha trasformato il paradosso della scelta in una condizione permanente, alimentata dall'illusione che infinite opzioni equivalgano a infinite possibilità di felicità. Pensate a Netflix: migliaia di titoli in decine di categorie, e un algoritmo che promette di conoscervi meglio di voi stessi. Eppure, quante volte avete scrollato per 20 o 30 minuti, per poi spegnere tutto o riguardare per la decima volta lo stesso comfort show?

Il problema non è l'algoritmo: è che la ricerca della scelta perfetta è diventata più importante dell'esperienza stessa. Barry Schwartz la chiama la «tirannia delle piccole decisioni»\footnote{Schwartz, B. (2004). \textit{The Paradox of Choice: Why More Is Less}. New York: Ecco.}: sprechiamo risorse cognitive preziose su scelte relativamente irrilevanti e arriviamo esausti a quelle che contano davvero.

\section{Maximizer e satisficer}

Herbert Simon, economista e psicologo premio Nobel, distinse due modi fondamentali di decidere\footnote{Simon, H. A. (1956). «Rational choice and the structure of the environment». \textit{Psychological Review}, 63(2), 129-138.}. Il \textbf{maximizer} cerca sempre la scelta ottimale: analizza tutte le opzioni, legge tutte le recensioni, confronta tutti i prezzi, e non si accontenta del «buono» perché vuole il «migliore». Il \textbf{satisficer} (neologismo da «satisfy» e «suffice») cerca invece una scelta «abbastanza buona»: stabilisce dei criteri minimi e si ferma alla prima opzione che li soddisfa. La Figura~\ref{fig:maximizer_satisficer} mette a confronto i due percorsi.

\begin{figure}[htbp]
	\centering
	\immagine{images/maximizer_satisficer.png}{Schema a due percorsi davanti a uno scaffale con molti prodotti. A sinistra il maximizer: una lunga freccia tortuosa che confronta ogni prodotto, termina in un acquisto con un'icona di dubbio e rimpianto. A destra il satisficer: una freccia breve che si ferma al primo prodotto che supera una soglia «abbastanza buono», termina con un'icona di soddisfazione.}
	\captionof{figure}{Due strategie di scelta: il maximizer esplora tutto e dubita, il satisficer si ferma al «buono abbastanza».}
	\label{fig:maximizer_satisficer}
\end{figure}

Diverse ricerche suggeriscono che i satisficer tendono a essere più felici delle loro scelte e della loro vita in generale. I maximizer, nonostante spesso facciano oggettivamente scelte «migliori», sono meno soddisfatti, più inclini al rimpianto, più vulnerabili alla depressione. Chi cerca la perfezione la trova raramente, e quando la trova non la riconosce; chi si accontenta del «buono abbastanza» finisce per sentirsi meglio.

\section{La FOMO e il fantasma delle vite non vissute}
\label{sec:p2c21-fomo}

L'era digitale ha aggiunto al paradosso una dimensione nuova: la FOMO (Fear Of Missing Out). Non solo dobbiamo scegliere tra infinite opzioni, ma i social media ci espongono di continuo alle scelte degli altri, e ogni nostra decisione viene confrontata non soltanto con le alternative scartate, ma con le vite apparentemente perfette che altri stanno vivendo. Instagram ci mostra le vacanze che non stiamo facendo, LinkedIn le carriere che non stiamo perseguendo, Facebook le relazioni che non abbiamo.

È una tortura cognitiva senza precedenti. I nostri antenati potevano invidiare il vicino di caverna, ma non erano bombardati 24/7 dalle vite curate di milioni di estranei.

\section{Il ritorno del vincolo}

C'è un'ironia profonda nel fatto che, in un'epoca di libertà senza precedenti, stiamo riscoprendo il valore dei vincoli: le app che bloccano i social media, i ristoranti con menu fisso, le capsule wardrobe che riducono l'armadio all'essenziale. Scegliamo consapevolmente di limitare le scelte, perché la vera libertà non sta nell'avere infinite opzioni, ma nell'avere lo spazio mentale per godere di quelle che abbiamo scelto.

È quello che Ulisse capì quando si fece legare all'albero della nave per resistere al canto delle sirene: a volte, per essere liberi, bisogna vincolarsi preventivamente. Ed è la lezione che l'arte ha sempre saputo. Un sonetto non è meno poetico perché deve rispettare 14 versi e uno schema di rime: spesso è proprio il vincolo formale a spingere alla creatività più alta.

\section{Perché il limite ci restituisce la scelta}

Se il problema nasce da un cervello costruito per poche opzioni, la via d'uscita non è allenarlo a reggerne infinite, ma riportare la scelta alla sua scala. Il banco delle sei marmellate non era una privazione: era un ambiente in cui la mente paleolitica poteva fare ciò per cui è fatta, confrontare poche cose concrete e decidere. Ridurre il campo prima di cominciare, che si tratti dei ristoranti del quartiere o di una lista già fatta, significa ricreare quel banco anche dentro un catalogo sterminato.

La seconda idea viene da Simon. Il satisficer non è più pigro del maximizer: ha semplicemente stabilito in anticipo che cosa gli basta, e così smette di cercare quando lo trova. Il rimpianto che tormenta il maximizer, e che la FOMO moltiplica, si nutre proprio di questa assenza di una soglia.

C'è infine il gesto di Ulisse, che decide una volta sola, a mente fredda, e poi si lascia legare. Ogni regola fissata in anticipo è una decisione sottratta alla fatica decisionale di domani, ed è per questo che il menu fisso, più che toglierci libertà, ci restituisce energia per le scelte che contano.

\bigskip

Dopotutto, anche in un buffet infinito, lo stomaco rimane della stessa dimensione. La domanda non è quante opzioni abbiamo, ma quanto siamo capaci di gustare quello che scegliamo. E per gustare davvero, a volte, bisogna smettere di assaggiare tutto.

C'è però un motivo se quell'assaggio infinito ci sembra così irrinunciabile: le vite che invidiamo e le scelte che crediamo giuste ci arrivano quasi sempre attraverso le storie di chi ce l'ha fatta. Di chi ha imboccato la stessa strada ed è rimasto indietro, invece, non racconta nessuno.
